\documentclass[12pt]{article}
\usepackage{amsmath}
\usepackage{amsfonts}
\usepackage{amssymb}
\usepackage{physics}
\usepackage{geometry}
\usepackage{parskip}
\geometry{a4paper, margin=0.7in}


\title{02YMECH - Solution 11}
\date{Assigned for the week of Dec 1, 2025}
\author{}

\begin{document}
\maketitle

\section*{Solutions}

\begin{enumerate}

\item Rod of length $l$ and mass $m$, with uniform linear density
\[
\lambda=\frac{m}{l}.
\]
Place the rod along the $x$-axis from $-l/2$ to $+l/2$. The observation point lies on the axis at distance $z$ from the center.

An element $dx$ contributes
\[
d\Phi = -\frac{G\lambda\,dx}{\sqrt{x^2+z^2}}.
\]
Thus,
\[
\Phi(z)=-G\lambda\int_{-l/2}^{l/2}\frac{dx}{\sqrt{x^2+z^2}}.
\]
Using
\[
\int \frac{dx}{\sqrt{x^2+a^2}}
=\ln\!\left(x+\sqrt{x^2+a^2}\right),
\]
we obtain
\[
\Phi(z)=-\frac{Gm}{l}
\ln\!\left(
\frac{\frac{l}{2}+\sqrt{z^2+\left(\frac{l}{2}\right)^2}}
{-\frac{l}{2}+\sqrt{z^2+\left(\frac{l}{2}\right)^2}}
\right).
\]

\item For a ring of radius $a$ and mass $m$, a point on the axis at distance $z$ from the center is at the same distance
\[
r=\sqrt{a^2+z^2}
\]
from every mass element. Hence,
\[
\Phi(z)=-G\int\frac{dm}{r}
=-\frac{Gm}{\sqrt{a^2+z^2}}.
\]

\item The total mass is
\[
M=\int_0^R 4\pi r^2\rho(r)\,dr
=\frac{4\pi\rho_0}{R}\int_0^R r^3dr
=\pi\rho_0 R^3,
\]
so that
\[
\rho_0=\frac{M}{\pi R^3}.
\]

The enclosed mass for $r<R$ is
\[
M(r)=\frac{4\pi\rho_0}{R}\int_0^r r'^3dr'
= M\frac{r^4}{R^4}.
\]

Therefore, the gravitational field is
\[
\mathbf g(r)=
\begin{cases}
-\displaystyle G M \frac{r^2}{R^4}\,\hat{\mathbf r},
& r<R,\\[1.2ex]
-\displaystyle \frac{GM}{r^2}\,\hat{\mathbf r},
& r\ge R.
\end{cases}
\]

\item For a uniform-density Earth of radius $R$ and mass $M$, the force on a particle at distance $r$ from the center is
\[
F=-\frac{GM(r)m}{r^2}
=-\frac{GMm}{R^3}\,r, ~~~ \text{where} ~~M(r)=M\frac{r^3}{R^3}.
\]
leading to the equation of motion
\[
\ddot r=-\frac{GM}{R^3}r.
\]
This is simple harmonic motion with
\[
\omega=\sqrt{\frac{GM}{R^3}}.
\]

The period is
\[
T=2\pi\sqrt{\frac{R^3}{GM}}
\approx 5.07\times10^3\ \text{s}
\approx 84.5\ \text{min}.
\]

\end{enumerate}



\end{document}
